\documentclass[11pt]{article}
\usepackage{fontspec}
\setmainfont{Times New Roman}
\setsansfont{Arial}
\setmonofont{Consolas}
\usepackage{amsmath,amssymb}
\usepackage{geometry}
\usepackage{microtype}
\geometry{a4paper,margin=3cm}
\setlength{\parindent}{1.25em}
\setlength{\parskip}{0pt}
\emergencystretch=10em
\allowdisplaybreaks
\providecommand{\Kfield}{\mathfrak{K}}
\providecommand{\Ssys}{\mathfrak{S}}
\providecommand{\Jdom}{\mathfrak{J}}
\providecommand{\Gdom}{\mathfrak{G}}
\providecommand{\Lsys}{\mathfrak{L}}
\providecommand{\Omegaint}{[\Omega]}
\providecommand{\eps}{\varepsilon}

\begin{document}

\section*{\S\ 15. Celočislene relativno cěle oblasti prvogo roda i regularne sistemy.}

Na město integralnoj bazy pri celočislenyh sistemah nastupa celočislena integralna baza, to jest konečno čislo celočislenyh polinomov sistema tako, že vsak celočisleny polinom sistema se može predstaviti kako cěla celočislena kombinacija sej konečnoj množiny.

Teoremy o celočislenej integralnoj bazi v sučnosti idut paralelno s predhodnymi, i samo dokazy trebujut male modifikacije.

Najprvo definicija $\mathrm{V}$ (\S\ 9) može se neposrědnje prěnesti.

{\itshape Definicija $\mathrm{V}'$: veličina $\xi$ nazyva se algebraično cěla i celočislena vzhodno k celočislenej relativno cělej oblasti $\Gdom_{n\rho}$, ako ona zadovolja někoj jednačbi, ktorej najvyšši koeficient jest jedinica iz $[\Omega]$, a ostale polinomy sut iz $\Gdom_{n\rho}$.}

Razširenje Gaussova teorema na polinomy s algebraično-celočislenymi koeficientami pokazuje imenno, kak v \S\ 9, že iz togo definicija može byti dobyta na ireducibilnoj jednačbi.

Iz definicije $\mathrm{V}'$ dalje slěduje prěnos definicije VII.

{\itshape Definicija $\mathrm{VII}'$: celočislena relativno cěla oblast $\Gdom_{n\rho}$ nazyva se oblastju prvogo roda, ako ona imaje kako oblast odobraženja relativno cělu oblast $\Gdom^*_{\rho\rho}$ tako, že vsak libovoljny celočisleny polinom od $x_1,\ldots,x_\rho$ algebraično cělo i celočisleno zavisi od $\Gdom^*_{\rho\rho}$.}

Da pokazati konečnost tyh oblastij prvogo roda, najprvo zamětimo, že kak v \S\ 10 iz definicije slěduje: najvyšši koeficient involucijskoj formy jest jedinica iz $[\Omega]$.

Po teoremu $\mathrm{VI}'$ (\S\ 14) zato vsak polinom $F(x)$ iz $\Gdom_{n\rho}$ dopušča predstavjenje:
\[
  F(x)=\frac{\Gamma(H_1(x),\ldots,H_\sigma(x))}{i}.
\]
Ako se tut na sistem, utvorjeny iz celočislenyh polinomov $\Gamma(H_1,\ldots,H_\sigma)$, priměni Hilbertov teorem II o celočislenej modulnoj bazi (Math. Ann. 36), izrečeny za cěle racionalne čislove koeficienty, v jego razširenju na algebraično-celočislene polinomy\footnote{Ako $w_1,\ldots,w_k$ označa fundamentalny sistem algebraično cělyh čisel iz $\Omega$, i ako postavimo
\[
\begin{gathered}
  \Gamma(H_1,\ldots,H_\sigma)
  =\Gamma_1(H)w_1+\cdots+\Gamma_k(H)w_k,
\end{gathered}
\]
togda $\Gamma_1(H),\ldots,\Gamma_k(H)$ imajut cěle racionalne čislove koeficienty. Priměnjenje teorema II na sistem, utvorjeny iz form $v_1\Gamma_1(H)+\cdots+v_k\Gamma_k(H)$, neposrědnje pokazuje važnost togo razširenja.}, togda iz togo slěduje obstojanje celočislenej integralnoj bazy, to jest:

{\itshape Teorem $\mathrm{VII}'$: vsaka celočislena relativno cěla oblast prvogo roda imaje celočislenu integralnu bazu.}

Dalje prěnosimo definiciju regularnyh sistemov:

{\itshape Definicija $\mathrm{VIII}'$: celočisleny sistem $\Ssys_{n\rho}$ nazyva se celočisleno-regularnym, ako v oblasti odobraženja $\Jdom^*_{\rho\rho}$ najmenšej soderžečej celočislenej integralnej oblasti obstojat $\rho$ polinomov tako, že rezultant členov najvyššej dimenzije jest jedinica iz $[\Omega]$, to jest $R_0=\eps$.}

Da prěnesti dokaz obstojanja integralnoj bazy, zamětimo, že pomočno tvrđenje iz \S\ 11 dopušča slědujuću ostrějšu formulaciju; ona slěduje iz jednačb (5) i (6) v \S\ 11 i iz okolnosti, že $A_1(\lambda),\ldots,A_\tau(\lambda)$ sut cěle celočislene funkcije od $\lambda_1,\ldots,\lambda_\rho$:

{\itshape Dostatno uslovje za to, že vsak libovoljny celočisleny polinom od $x_1,\ldots,x_\rho$ algebraično cělo i celočisleno zavisi od $\rho$ celočislenyh polinomov, jest to, že rezultant členov najvyššej dimenzije tyh polinomov jest jedinica iz $[\Omega]$, to jest $R_0=\eps$.}

Iz sej pomoćnoj teoremy slěduje točno kak v \S\ 12 za vsak polinom $F(x)$ celočisleno-regularnogo sistema predstavjenje (3) iz \S\ 12, gdě nyně $\Gamma_1(G_i),\ldots,\Gamma_k(G_i)$ sut celočislene polinomy od $G_i$. I dalje, kak v \S\ 12, priměnjenje Hilbertova teorema II (Math. Ann. 36) v jego razširenju na algebraično-celočislene polinomy, vměstě s definiciju najmenšej soderžečej celočislenej integralnej oblasti, davaje obstojanje celočislenej integralnoj bazy; to jest:

{\itshape Teorem $\mathrm{IX}'$: vsak celočisleno-regularny sistem imaje celočislenu integralnu bazu.}

Kak primjer celočislenej relativno cělej oblasti prvogo roda, v analogiji s primjerom 1 v \S\ 13, ukazujemo na cělost $\Gdom_{nn}$ celočislenyh polinomov Lagrangevoj rodovoj oblasti. Že $\Gdom_{nn}$ jest celočislena oblast prvogo roda, slěduje iz identičnosti:
\[
  x_k^n-\sigma_1(x)x_k^{n-1}+\cdots\pm\sigma_n(x)=0
  \qquad(k=1,2,\ldots,n);
\]
poněže rezultant elementarnyh simetričnyh funkcij imaje cěnnost $\pm1$, $\Gdom_{nn}$ jest takože celočisleny regularny sistem.

Dalejši primjer jest, po pomoćnom teoremu iz \S\ 14, dany cělostju cělyh celočislenyh simetričnyh funkcij od $n$ redov veličin.

Erlangen, maj 1914.

\end{document}
